\paragraph{Related arithmetic methods.}
The elementary denominator argument has a direct precedent in
W.~Zudilin, \emph{Well-poised hypergeometric service for Diophantine problems
of zeta values}, Journal de Th\'eorie des Nombres de Bordeaux
15 (2003), 593--626, especially equations (14)--(16):
\url{https://www.numdam.org/item/JTNB_2003__15_2_593_0.pdf}.
There the basic ingredients are integer-valued binomial polynomials,
their derivatives after lcm multiplication, and factorial-normalized
reciprocal blocks with integer residues. These ingredients justify the
baseline arithmetic above; they do not automatically supply the stronger
common divisor needed for this particular trajectory.

C.~Krattenthaler and T.~Rivoal's \emph{Hyperg\'eom\'etrie et fonction z\^eta de
Riemann}, Memoirs of the AMS 186 (2007), no.~875,
\url{https://arxiv.org/abs/math/0311114}, develops stronger denominator
results using hypergeometric identities. Its general theme is relevant,
but no theorem from that paper is being asserted to apply to this exact
parameter set without checking its hypotheses.

For current qualitative context, Li Lai's preprint
\emph{The linear independence of $1$, $\zeta(3)$, and $\zeta(4)$},
submitted 5 October 2026,
\url{https://arxiv.org/abs/2610.06838}, states precisely that linear
independence result. Lemma 2.2 uses double- and simple-pole coefficients
and harmonic terms of orders two through four, closely matching the
readout structure here. Its method uses determinants and residue-class
arithmetic. We have checked the statement and relevant formulas, not
independently audited its full proof; this work makes no novelty claim
for the qualitative irrationality of every fixed mixed combination.
